\section{The compact Hall atlas obstruction}
\label{sec:compact-hall-atlas-obstruction}

\providecommand{\Z}{\mathbb Z}
\providecommand{\C}{\mathbb C}
\providecommand{\GammaBPS}{\Gamma_{\mathrm{BPS}}}
\providecommand{\Gammaeff}{\Gamma_{\mathrm{eff}}}
\providecommand{\Gammaact}{\Gamma_{\mathrm{act}}}
\providecommand{\Hall}{\mathcal H}
\providecommand{\Perf}{\operatorname{Perf}}
\providecommand{\Coh}{\operatorname{Coh}}
\providecommand{\Hilb}{\operatorname{Hilb}}
\providecommand{\Ext}{\operatorname{Ext}}
\providecommand{\RHom}{\operatorname{RHom}}
\providecommand{\BM}{\operatorname{BM}}
\providecommand{\Crit}{\operatorname{Crit}}
\providecommand{\Prim}{\operatorname{Prim}}
\providecommand{\sdim}{\operatorname{sdim}}
\providecommand{\BPS}{\mathrm{BPS}}

The object under attack is the compact oriented critical Hall atlas
postulated in the compact Hall factorization note.  The
question is whether it can be constructed directly for
\(X=S\times E\), \(S\) a K3 surface, from the curve-counting geometry
which gives the Oberdieck--Pixton scalar square.

The answer is negative at the present level of input.  One can write the
candidate moduli problem and the conditional Hall atlas.  One cannot yet
descend the reduced \(E\)-quotient scalar theory to a Hall convolution
atlas whose protected primitive bracket has denominator
\(\mathfrak g_{\Delta_5}\).

\subsection{The candidate scalar moduli problem}

Let \(\beta_h\in H_2(S,\Z)\) be primitive with
\(\langle \beta_h,\beta_h\rangle=2h-2\).  The Hilbert-scheme DT branch
uses the projective moduli space
\[
  \Hilb^{\ell}\bigl(X,(\beta_h,d)\bigr)
\]
of one-dimensional subschemes \(Z\subset X\) with
\([Z]=(\beta_h,d)\) and \(\chi(\mathcal O_Z)=\ell\), together with the
translation action of \(E\) on the second factor of \(X=S\times E\).
The numerical invariant displayed in the manuscript is
\[
  \mathbf{DT}^{X}_{\ell,(\beta_h,d)}
  =
  \int_{\Hilb^{\ell}(X,(\beta_h,d))/E}\nu\,de ,
\]
where \(\nu\) is Behrend's constructible function
\cite{BEHREND}, \cite{DT}.  In the OP variables,
\[
  q=q_{\mathrm{OP}}=t_{\mathrm{DT}},\qquad
  r=p_{\mathrm{OP}}=p_{\mathrm{DT}},\qquad
  s=\widetilde q_{\mathrm{OP}}=q_{\mathrm{DT}},
\]
and
\[
  \chi_{10}^{\mathrm{OP}}
  =
  4096^{-1}\Delta_5^2,\qquad
  Z^X_{\mathrm{OP}}=-4096\,\Delta_5^{-2}
\]
\cite{OB}, \cite{OPand}.  The Igusa product uses the charge
\(\gamma=(n,l,m)\) in
\[
  q^n r^l s^m,\qquad
  \alpha(n,l,m)=2n f_2-l f_3+2m f_{-2}.
\]

\begin{proposition}[Scalar critical geometry]
\label{prop:scalar-critical-geometry-only}
The Hilbert-scheme branch supplies a proper scalar DT moduli problem,
with symmetric obstruction theory and Behrend weighting.  It does not,
by itself, supply a Hall atlas.
\end{proposition}

\begin{proof}
The Hilbert scheme is the moduli space of ideal sheaves
\(I_Z\subset\mathcal O_X\) with fixed curve and Euler data.  It is
proper because \(X\) is projective.  The Calabi--Yau threefold
obstruction theory is symmetric, and the Behrend function gives the
Euler realization of the local critical structure.

Hall convolution needs an extension-closed moduli problem.  The class of
ideal sheaves, or stable pairs in a fixed chamber, is not closed under
arbitrary extensions in \(\Coh(X)\).  An exact sequence
\[
  0\longrightarrow I_{Z_1}\longrightarrow F\longrightarrow I_{Z_2}
  \longrightarrow 0
\]
need not have \(F\) an ideal sheaf.  Thus the Hilbert scheme is a scalar
DT coefficient space, not the Hall stack on which the primitive
commutator is defined.  A Hall atlas must be built on an abelian or
derived heart containing these objects and closed under extensions.
\end{proof}

\subsection{The raw finite-truncation obstruction}

The previous Hall note used finite lower charge sets
\(\Gamma_{N,R}\subset\Gammaeff\).  This is correct only after a stability
sector and Harder--Narasimhan finiteness have been supplied.  It is false
as a statement about the raw Igusa cone.

\begin{lemma}[No raw finite lower sets]
\label{lem:no-raw-finite-lower-sets}
Let \(L\subset\Gammaeff\) have the lower-set property
\[
  \gamma+\eta\in L,\qquad \gamma,\eta\in\Gammaeff
  \quad\Longrightarrow\quad
  \gamma,\eta\in L .
\]
If \(L\) contains any charge \((n,l,m)\) with \(n>0\) and \(m>0\), then
\(L\) is infinite.
\end{lemma}

\begin{proof}
For every \(a\in\Z\),
\[
  (n,l,m)=(n,a,0)+(0,l-a,m).
\]
Both summands lie in \(\Gammaeff\): the first has \(m=0,n>0\), and the
second has \(m>0,n=0\).  Hence a lower set containing \((n,l,m)\) must
contain \((n,a,0)\) and \((0,l-a,m)\) for all \(a\in\Z\).  It is
therefore infinite.
\end{proof}

\begin{corollary}[Finite-first means HN-first]
\label{cor:finite-first-means-hn-first}
The finite-first Hall construction cannot be indexed by finite lower
subsets of the raw cone \(\Gammaeff\).  It must be indexed by finite
Harder--Narasimhan sector truncations
\[
  \Gamma_{\sigma,S}(N,R)
  =
  \{\gamma\in\Gamma^{ss}_{\sigma,S}\mid
  \|\gamma\|\le N,\ |Z_\sigma(\gamma)|\le R\},
\]
enlarged by the actual HN factors of objects in the truncation.  The
finiteness of these sets is a consequence of a stability condition,
boundedness, and the support property.  It is not a consequence of the
Borcherds product chamber alone.
\end{corollary}

\begin{proof}
Lemma~\ref{lem:no-raw-finite-lower-sets} rules out raw finite lower
sets.  In a genuine Hall construction, only decompositions which occur
as HN factors in a strict stability sector must be included.  The
support property bounds the charge norm in bounded central-charge
radius, and boundedness of semistable moduli makes the truncation
finite.  These are stability-theoretic hypotheses, not modular-form
data.
\end{proof}

\begin{remark}[Active support is not a Hall truncation]
The active denominator support
\[
  \Gammaact=\{\gamma=(n,l,m)\in\Gammaeff\mid f(nm,l)\ne0\}
\]
is finite in the \(l\)-direction after fixing \(n,m\), because
\(\phi_{0,1}\) is weak of index one and
\[
  f(nm,l)\ne0\quad\Longrightarrow\quad 4nm-l^2\ge -1.
\]
This controls the Borcherds cone.  It does not control the Hall stack:
inactive charges may occur as extensions and may be needed before
primitive extraction kills their protected index.
\end{remark}

\subsection{The orientation datum}

Let \(\mathfrak M\) be a derived moduli stack of objects in a chosen
heart inside \(\Perf(X)\).  The Calabi--Yau structure gives a
\((-1)\)-shifted symplectic form on \(\mathfrak M\).  Locally on the
classical truncation this gives critical charts
\[
  \mathfrak M_\gamma|_U\simeq [\Crit(f_{\gamma,U})/G_{\gamma,U}].
\]
The determinant line is
\[
  \mathcal L_\gamma
  =
  \operatorname{sdet}\RHom_X(\mathcal E_\gamma,\mathcal E_\gamma),
\]
for the universal object when it exists, or its descent on the moduli
stack otherwise.  A strong orientation is a square root
\[
  \mathscr L_{o,\gamma}^{\otimes2}\simeq\mathcal L_\gamma
\]
together with exact-triangle coherences on the extension stacks.

\begin{proposition}[Orientation needed for Igusa]
\label{prop:orientation-needed-for-igusa}
An oriented critical atlas for \(K3\times E\) requires more than the
existence of local square roots.  It requires an \(E\)-equivariant strong
orientation compatible with Hall convolution, reduced \(E\)-quotient
integration, and the automorphic character
\[
  \epsilon_o=\nu_{\Delta_5}.
\]
\end{proposition}

\begin{proof}
The local square root orients the vanishing-cycle coefficient.  Hall
multiplication also uses Thom--Sebastiani transport on the stack of
extensions, hence needs exact-triangle compatibility of the square
roots.  The OP scalar theory then divides by the translation action of
\(E\), so the orientation must be equivariant for that action.  Finally,
the Borcherds denominator carries the character \(\nu_{\Delta_5}\).  No
local square-root theorem identifies its induced boundary character with
\(\nu_{\Delta_5}\); this is an additional comparison datum.
\end{proof}

\subsection{The reduced quotient obstruction}

The \(E\)-quotient is the geometric operation which makes the OP scalar
nonzero.  It is also the point at which Hall convolution stops being
automatic.

\begin{proposition}[The reduced quotient is not automatically Hall]
\label{prop:reduced-quotient-not-hall}
Let \(\mathfrak M_\gamma\) be an \(E\)-equivariant moduli stack and let
\[
  \mathfrak E_{\gamma,\eta}
  \longrightarrow
  \mathfrak M_{\gamma+\eta}
\]
be the extension stack defining Hall multiplication before reduction.
The scalar quotient stacks \(\mathfrak M_\gamma/E\) do not inherit a
Hall product unless one constructs a relative-position or equivariant
descent map comparing the diagonal \(E\)-action on extensions with the
two independent \(E\)-actions on the input factors.
\end{proposition}

\begin{proof}
The unreduced Hall correspondence is equivariant for simultaneous
translation of the subobject, quotient object, and middle object.  Thus
it sees the diagonal \(E\)-action.  The product of reduced scalar
quotients,
\[
  (\mathfrak M_\gamma/E)\times(\mathfrak M_\eta/E),
\]
has forgotten the relative \(E\)-position of the two input objects.  But
extensions depend on this relative position.  Therefore the convolution
map on reduced quotient spaces is not obtained by formal descent from
the unreduced Hall correspondence.  One must either keep the
\(E\)-equivariant Hall algebra until after convolution, or construct an
explicit reduced convolution kernel which retains the relative-position
parameter.  The manuscript has neither construction.
\end{proof}

\subsection{Conditional construction}

\begin{definition}[Igusa compact Hall atlas datum]
\label{def:igusa-compact-hall-atlas-datum}
An Igusa compact Hall atlas datum for \(X=S\times E\) consists of:
\begin{enumerate}
\item an extension-closed heart \(\mathcal A_\sigma\subset\Perf(X)\)
  with a stability condition \(\sigma\), HN property, support property,
  and bounded semistable moduli;
\item a charge map
  \[
    \operatorname{cl}\colon K_0(\mathcal A_\sigma)\longrightarrow
    \GammaBPS=\Z^3
  \]
  whose Igusa chamber is \(\Gammaeff\) and whose Borcherds degree is
  \(\alpha(n,l,m)=2nf_2-lf_3+2mf_{-2}\);
\item for each strict sector \(S\), finite HN truncations
  \(\Gamma_{\sigma,S}(N,R)\) as in
  Corollary~\ref{cor:finite-first-means-hn-first};
\item derived Artin substacks
  \(\mathfrak M_{\sigma,S,\gamma}\) of \(\sigma\)-semistable objects,
  with local critical presentations and vanishing-cycle coefficients;
\item an \(E\)-equivariant strong orientation
  \(o=\{\mathscr L_{o,\gamma}\}\) compatible with exact triangles;
\item extension correspondences
  \[
    \mathfrak M_{\gamma}\times\mathfrak M_{\eta}
    \xleftarrow{\ p\ }
    \mathfrak E_{\gamma,\eta}
    \xrightarrow{\ q\ }
    \mathfrak M_{\gamma+\eta}
  \]
  for the chosen sector, with compact-support pull-push and
  Thom--Sebastiani transport;
\item Cech/Ran descent for the local critical charts and HN sector
  descent for ordered sector decompositions;
\item an \(E\)-equivariant or reduced integration map compatible with
  Hall convolution and inducing the OP scalar normalization.
\end{enumerate}
\end{definition}

\begin{theorem}[Conditional compact Hall atlas]
\label{thm:conditional-compact-hall-atlas}
If an Igusa compact Hall atlas datum in the sense of
Definition~\ref{def:igusa-compact-hall-atlas-datum} is supplied, then
for each \((N,R)\) there is a finite critical Hall truncation
\[
  \mathsf F_{N,R}\Hall^{\mathrm{or}}_{X,\sigma,S,o}
  =
  \bigoplus_{\gamma\in\Gamma_{\sigma,S}(N,R)}
  H^{\BM}_{E,G_\gamma}
  \bigl(\mathfrak M_{\sigma,S,\gamma},
  \varphi_{\gamma}\otimes\mathscr L_{o,\gamma}\bigr),
\]
with product defined by Hall pull-push and oriented
Thom--Sebastiani transport.  These truncations form a completed
sectorial Hall factorization cosheaf
\[
  \widehat{\Hall}^{\mathrm{or}}_{X,\sigma,S,o}
  =
  \varprojlim_{N,R}
  \mathsf F_{N,R}\Hall^{\mathrm{or}}_{X,\sigma,S,o}.
\]
\end{theorem}

\begin{proof}
All operations are finite at fixed \((N,R)\) by the HN truncation and
boundedness hypotheses.  The local critical presentations and strong
orientation glue the vanishing-cycle coefficients.  The extension
correspondences define the Hall pull-push.  Thom--Sebastiani and the
orientation coherences identify the coefficient systems under
extensions.  Associativity is the equality of the two projections from
the stack of two-step filtrations.  Cech/Ran descent and HN uniqueness
assemble the finite products over strict sector decompositions.  Passing
to the inverse limit gives the completed cosheaf.
\end{proof}

\subsection{Failure theorem}

\begin{theorem}[What is not yet constructed for \(K3\times E\)]
\label{thm:k3e-atlas-not-yet-constructed}
The present manuscript, together with the current local agent material
and Vol.~III anchors, does not construct the Igusa compact Hall atlas
datum of Definition~\ref{def:igusa-compact-hall-atlas-datum}.
\end{theorem}

\begin{proof}
The missing points are independent.

First, the OP theorem constructs a scalar reduced curve-counting series.
It does not choose an extension-closed heart whose semistable objects are
the objects counted by that scalar theory.  The Hilbert-scheme spaces
are proper coefficient spaces, but Proposition~\ref{prop:scalar-critical-geometry-only}
shows that they are not themselves the Hall stack.

Second, finite lower truncations cannot be read from
\(\Gammaeff\) alone.  Lemma~\ref{lem:no-raw-finite-lower-sets} proves
that raw lower sets containing an interior charge are infinite.  Hence
the finite-first construction requires a stability condition with HN
finiteness and support control.

Third, the reduced \(E\)-quotient used in the scalar invariant does not
automatically preserve Hall convolution.  Proposition~\ref{prop:reduced-quotient-not-hall}
identifies the missing relative-position or equivariant descent map.

Fourth, the strong orientation required for vanishing cycles is not the
same statement as the automorphic character comparison
\(\epsilon_o=\nu_{\Delta_5}\).  The latter is a boundary comparison with
the Borcherds denominator character.

Fifth, the Vol.~III compact-CY3 Hall construction supplies a conditional
template over a realized stability chamber, strong orientation, and
admissible motivic coefficients.  It does not, in the local sources
read here, prove that the \(K3\times E\) OP chamber supplies those
inputs or that the resulting protected primitives have signed
dimensions \(f(nm,l)\).

Therefore the constructed object remains conditional.  The scalar square
and the denominator algebra determine the expected boundary character;
they do not construct the compact oriented critical Hall atlas.
\end{proof}

\subsection{The missing maps}

The construction closes only after the following maps are written and
checked.
\[
\begin{gathered}
  \mathrm{HN}_{\sigma,S}:
  \mathfrak M_{\sigma,S}
  \longrightarrow
  \mathfrak M_{\sigma,S_1}\widetilde{\times}\cdots
  \widetilde{\times}\mathfrak M_{\sigma,S_k},\\
  \mathrm{TS}_{o,\gamma,\eta}:
  \varphi_{\gamma}\boxtimes\varphi_{\eta}\otimes
  \mathscr L_{o,\gamma}\boxtimes\mathscr L_{o,\eta}
  \longrightarrow
  \varphi_{\gamma+\eta}\otimes\mathscr L_{o,\gamma+\eta},\\
  \mu^{E}_{\gamma,\eta}:
  H^{\BM}_{E}(\mathfrak M_\gamma,\varphi_\gamma\otimes\mathscr L_{o,\gamma})
  \otimes_{H^\bullet(BE)}
  H^{\BM}_{E}(\mathfrak M_\eta,\varphi_\eta\otimes\mathscr L_{o,\eta})
  \longrightarrow
  H^{\BM}_{E}(\mathfrak M_{\gamma+\eta},
  \varphi_{\gamma+\eta}\otimes\mathscr L_{o,\gamma+\eta}),\\
  \mathrm{Red}_E:
  H^{\BM}_{E}(\mathfrak M_\gamma,\varphi_\gamma\otimes\mathscr L_{o,\gamma})
  \longrightarrow
  H^{\BM}(\mathfrak M_\gamma/E,
  \varphi_\gamma^{\mathrm{red}}\otimes\mathscr L_{o,\gamma}^{\mathrm{red}}),\\
  \chi^{\mathrm{prot}}:
  K_0\Prim
  \widehat{\Hall}^{\mathrm{or}}_{X,\sigma,S,o}
  \longrightarrow
  \widehat{\mathbb T}_{\Gammaeff},
  \qquad
  \sdim\Prim_\gamma=f(nm,l),\\
  \epsilon_o:\{\hbox{strong orientations}\}\longrightarrow
  \{\hbox{Borcherds characters}\},
  \qquad
  \epsilon_o(o)=\nu_{\Delta_5}.
\end{gathered}
\]
The first two maps are local-to-global atlas maps.  The third and
fourth are the equivariant/reduced Hall maps.  The fifth is the
protected primitive integration.  The last is the automorphic boundary
orientation comparison.

\subsection{Dirac-standard rewrite block}

The manuscript presentation is:

\begin{quote}
Fix \(X=S\times E\).  The OP theorem computes a reduced scalar series on
the \(E\)-quotient of the curve-counting moduli.  A Hall algebra is not
formed on that quotient.  It is formed before reduction, on a
charge-graded oriented critical stack of objects closed under
extensions.

Thus the missing object is precise.  Choose a stability condition
\(\sigma\) on an extension-closed heart in \(\Perf(X)\), with HN
finiteness and support property.  Choose an \(E\)-equivariant strong
orientation of the induced \((-1)\)-shifted critical stack.  Construct
the extension correspondences and their Thom--Sebastiani orientation
transport.  Only then take the finite HN completion and the reduced
integration.

If the protected primitive integration gives
\[
  \sdim\Prim_\gamma=f(nm,l),
  \qquad \gamma=(n,l,m)\in\Gammaeff,
\]
and if the orientation character is \(\nu_{\Delta_5}\), the protected
determinant is \(64^{-1}\Delta_5(2Z)\).  The bracket is still a further
claim: it requires the Hall commutator to satisfy the
Gritsenko--Nikulin Borcherds relations.
\end{quote}

This is the whole physics at this point.  The scalar theory is the
measured shadow.  The Hall atlas is the mechanism.  The denominator
algebra is obtained from the mechanism only after the primitive and
bracket comparisons are proved.
